We have analysed a compressible hyperelastic corner under mixed Dirichlet--Neumann boundary conditions and shown that
the conditioning collapse first observed numerically is, in its essence, an analytical phenomenon: a corner-localised
\emph{creasing} bifurcation seated at the geometric junction, of which the divergence of the linearised tangent is a
symptom rather than the substance.

The central finding is a dichotomy in the asymptotic behaviour of the linearised tangent operator, governed by the
loading direction and the volumetric penalty. In tension the logarithmic barrier $f_{\log}$ yields a \emph{benign}
singularity: the tangent stays strongly elliptic and Fredholm, the standard Kondratiev edge exponents survive as the
intermediate asymptotics outside the nonlinear crossover scale --- at the vertex itself the extreme stretch screens the
volumetric coupling and the tangent stabilises to a decoupled Laplacian with no singular exponent at all --- and the
well-posedness of the linearised step holds in the weighted corner spaces (Theorem~\ref{thm:A}). Under compression the
local volume contraction drives an algebraic blow-up of the operator norm. We proved that this breakdown is a
\emph{structural bifurcation}, not a material loss of ellipticity (Lemma~\ref{lem:SEcompress}): the corner remains
pointwise strongly elliptic throughout, while a localised surface instability preempts pointwise collapse. The apparent
tangent catastrophe is, on every physical branch, regularised by the kinematic relief of the buckled mode
(Assumption~\ref{ass:rot}, Lemma~\ref{lem:relief}); a load-controlled body snaps to a finite-amplitude crease at a load
reduced by only $\mathcal{O}(\epsilon^{2/3})$, and only an algorithm that suppresses the snap and tracks the unstable
continuation sees the $\epsilon^{-4/3}$ conditioning degradation, confined to a tube of vanishing area and integrable in
$L^q$ for every $q<\tfrac12$ (Theorem~\ref{thm:reg}, Proposition~\ref{prop:narrow}).

A tangent-cone blow-up reduction localises the entire phenomenon to a \emph{universal} model wedge, governed only by the
Poisson ratio $\nu$ and a single normalised stress-intensity amplitude $\hat K$ (Proposition~\ref{prop:blowup}). Against
this backdrop the imperfect-bifurcation analysis is sharp on two counts and honest about a third. The quadratic energy
coefficient vanishes \emph{exactly} by reflection symmetry, with no Lyapunov--Schmidt correction (the parity argument of
\S\ref{sub:symmetry}). The subcritical cubic coefficient is tip-convergent yet does \emph{not} localise: it retains an
explicit dependence on the global stress-intensity factor, a direct consequence of the characteristic edge exponent
satisfying $\alpha>\tfrac12$. The two surviving hypotheses --- buckling-precedes-collapse (BC) and the subcritical sign
of $b$ --- sit on opposite sides of this blow-up dichotomy (Remark~\ref{rem:bclocalises}): the former compares
amplitude-free, $\nu$-only \emph{stretch} thresholds and is settled by a single wedge computation, while the latter is an
\emph{integral} whose sign is an $\mathcal{O}(1)$ geometric competition decided per configuration. The scale-covariant
Mellin--Chebyshev spectral solver resolves the singular exponents and buckling thresholds with exponential accuracy and
supplies the geometry-dependent diagnostics --- the surface-instability stretch $\lambda_\star(\nu)$, the rotation
defect, and the sign of $b$ --- that the analysis isolates but cannot decide in closed form.

What this concrete realisation buys is a sharp reduction of hypotheses. Fixing the clamped--free neo-Hookean material and
geometry discharges, one by one, the spectral postulates of the abstract theory into theorems (the ledger of
Table~\ref{tab:status}): the material bounds become explicit, the critical load is constructed from a localised surface
instability (Lemma~\ref{lem:biot}), mode attainment is \emph{proved} under (BC) (Lemma~\ref{lem:attain}), simplicity is
reduced from an outright postulate to genericity within the symmetry sector carrying the mode, the quadratic coefficient
vanishes by symmetry (Lemma~\ref{lem:volcoeff}), and the exact first-order volume preservation once posited by the relief
is weakened to the asymptotic-rotation hypothesis, automatic at the bifurcation (Lemma~\ref{lem:rotreg}) and numerically
supported toward collapse. What is left after this reduction is not a residue of opaque spectral assumptions but exactly
two \emph{concrete}, numerically checkable conditions --- buckling precedes pointwise collapse, and the reduced cubic is
subcritical. Crucially, these are the \emph{same} two conditions that the classical creasing problem on a flat surface
also leaves to computation \cite{Biot1963,Hong2009}: they are inherited here, not manufactured by the corner geometry,
which adds the singular pre-stress and the localisation but no new open hypothesis.

Both surviving conditions are, in practice, decided by the spectral computation of \S\ref{sec:numerical}. The sign of
the cubic coefficient is the more complete of the two: the graded-cell solve of \S\ref{sub:numb} returns
$b<0$ throughout the tested $(\nu,\hat K)$ rectangle --- the normalisation-invariant feedback ratio stays below
its critical value $-1$ at every tested resolution, with the smallest margin near the incompressible corner, and the
inversion-free $\Gamma$ certificate of Corollary~\ref{cor:cert} now confirms the sign at every rectangle point
without the ill-conditioned reduced solve --- so the cubic clause of subcriticality is numerically supported and
internally cross-validated, though on a surrogate cell whose refinement behaviour is not yet settled. The quintic has
now been computed as well, through the sixth-order reduction of Lemma~\ref{lem:quintic}, and the verdict runs the
other way: $d<0$ throughout, so the polynomial normal form does not saturate --- the corner inherits the strongly
nonlinear, finite-amplitude character of the crease, and the arrested post-snap state rests on the non-degeneracy
hypothesis (ND) rather than on a quintic fold. What remains genuinely beyond reach here is
analytic rather than numerical: a \emph{closed-form, configuration-free} criterion for $\operatorname{sign} b$, which the
homogeneity count of Proposition~\ref{prop:bcount} shows cannot exist in the localised form available for the stretch
thresholds, precisely because $\alpha>\tfrac12$ keeps the cubic integral globally $\hat K$-dependent. The loading side of
(BC) is likewise computed: \S\ref{sub:numbiot} delivers the amplitude-free stretch threshold $\lambda_\star(\nu)$ to
spectral accuracy. The one item left open as a theorem is the inner vertex theory for the logarithmic barrier
(Remark~\ref{rem:core}): away from the tip the model-wedge operator and its corner pencil are pinned to
constant-coefficient linear elasticity, but a quantitative, uniform-to-the-tip spectral estimate --- and with it the
lower bound on the barrier-regularised collapse load $t_{\mathrm{coll}}$ that would turn the computed stretch ordering
into the load ordering $t_c<t_{\mathrm{coll}}$ --- is the irreducible analytic core of (BC). Beyond this, the natural
continuations are the finite-amplitude creased state past the snap, the extension from the model wedge to fully
three-dimensional edges with curvature, and the analogous dichotomy for other volumetric penalties, where the same
competition between barrier growth and volumetric contraction should organise the asymptotics.
